% !TEX root = main.tex
%
% Tate-conilpotent bar-cobar Quillen equivalence.
%
% Relative model-categorical promotion of the cochain-level Tate
% bar-cobar identity of \ref{lem:continuous-bar-cobar}.  The Quillen
% statements below live in the admissible bounded-below
% Mittag-Leffler Tate envelope fixed in
% \ref{stmt:tate-model-envelope}.  They do not assert a model
% structure on every raw Tate vector space and do not solve the
% analytic Casimir / heat-kernel residual.
%
% We use
% Loday-Vallette \emph{Algebraic Operads} (2012) Ch. 11-13 sign and
% suspension conventions. The Tate categorical machinery is
% Costello-Gwilliam Vol. I \S\S 6.1, 6.4, 7.2 (factorization
% algebra, local constancy, extension-from-a-basis vocabulary)
% combined with Hinich \emph{Homological algebra of homotopy
% algebras} (1997) \S 2 transfer theorem.
%
% Sign conventions: Koszul. Suspension
% $s\colon V\to V[1]$ shifts homological degree up by one. The bar
% and cobar functors satisfy
%   $\mathrm{Bar}(A)=T^c(s\overline A),\quad
%    \Omega(C)=T(s^{-1}\overline C)$,
% with differentials assembled from $d_A,m_A,\Delta_C,d_C$ in the
% standard Loday-Vallette sign (\S 2.2.6, \S 11.1).

\subsection[Tate Quillen equivalence and the infinity-categorical upgrade]{Tate Quillen equivalence and the $\infty$-categorical
upgrade of the cochain bar-cobar identity}\label{sec:tate-quillen}

\subsubsection{Categorical home}

Let $\Cat{TateVec}$ be the closed symmetric monoidal category of
complete filtered Tate $\C$-vector spaces of \alpharef. The
filtration on every object is taken \emph{complete and Hausdorff}:
$V=\varprojlim_w V_{\leq w}$ and $\bigcap_w F^w V=0$. Concrete
morphisms are filtration-preserving continuous chain maps. The
working ambient category is the bounded-below filtered subcategory
$\Cat{TateChain}_{\geq 0}$ of non-negatively filtered chain
complexes in $\Cat{TateVec}$, large enough to contain every CE
  chain coalgebra and completed enveloping algebra used here.  The
  model-categorical statements below are made in a fixed admissible
  envelope, not in the raw category of all Tate spaces.

\begin{stmt}[Admissible Tate model envelope]\label{stmt:tate-model-envelope}
  Fix a bounded-below, locally presentable symmetric monoidal model
  envelope $\Cat{TateChain}_{\geq0}^{\mathrm{adm}}$ containing the
  CE chains, completed enveloping algebras, completed symmetric
  algebras, and finite-window inverse systems used in this
  manuscript.  Weak equivalences are admissible filtered
  quasi-isomorphisms: quasi-isomorphisms on every finite quotient
  and on associated graded.  The tensor product is the completed
  Tate tensor product on Mittag-Leffler systems, and filtered inverse
  limits with surjective transition maps are exact.  All Quillen
  statements below are relative to this envelope.
\end{stmt}

\begin{defn}[Tate cofibrantly generated model structure in the envelope]
  \label{defn:tate-model}
  Define generating cofibrations $I$ and acyclic cofibrations $J$
  by
  $$I=\bigl\{F^w S^{n-1}\hookrightarrow F^w D^n
       \mid n\in\Z_{\geq 0},\ w\in\N\bigr\},\qquad
    J=\bigl\{0\hookrightarrow F^w D^n
       \mid n\in\Z_{\geq 0},\ w\in\N\bigr\},$$
  where $S^{n-1},D^n$ are the standard sphere and disk chain
  complexes in $\Cat{Vec}$ and the Tate window $F^w(-)$ takes the
  $w$-th piece of the filtration. A morphism $f\colon V\to W$ in
  $\Cat{TateChain}_{\geq 0}^{\mathrm{adm}}$ is:
  \begin{enumerate}
    \item a \emph{weak equivalence} if it is an admissible filtered
      quasi-isomorphism in the sense of \alpharef:
      a quasi-isomorphism on every quotient $V/F^wV\to W/F^wW$ and
      on the associated graded;
    \item a \emph{cofibration} if it is in the class generated by
      $I$ under transfinite composition, retracts, and pushouts;
    \item a \emph{fibration} if it has the right lifting property
      against $J$.
  \end{enumerate}
\end{defn}

\begin{thm}[Existence inside the admissible envelope]
  \label{thm:tate-model-structure}
  Under Statement~\ref{stmt:tate-model-envelope},
  the triple
  \[
    \bigl(\Cat{TateChain}_{\geq 0}^{\mathrm{adm}},I,J\bigr)
  \]
  is a cofibrantly generated model category. The class of weak equivalences is
  closed under retracts and satisfies the two-out-of-three
  property; cofibrations and fibrations form complementary lifting
  classes; the small-object argument terminates after countably
  many steps because every Tate window is finite-dimensional and
  every generating cofibration changes only one window.
\end{thm}

\begin{proof}
  Quillen's small-object argument (\cite{quillen-htpy}*{\S II.3})
  applies in the admissible envelope once the following are
  verified:
  \begin{enumerate}[label=(\roman*)]
    \item $I$ and $J$ permit the small-object argument:
      the domains and codomains of the generating maps are small
      with respect to relative $I$-cell complexes because each
      generator lives in one finite-dimensional Tate window
      $V_{\leq w}$, hence is $\aleph_0$-small in the admissible
      envelope.
    \item Every map in $J$-cell is a weak equivalence: pushouts of
      acyclic generating cofibrations are levelwise acyclic by
      finite-dimensional standard chain-complex theory in each
      window, and the filtered colimit of windowwise acyclic
      cofibrations is acyclic by Mittag-Leffler exactness in the
      Tate filtration (\alpharef\ item~1).
    \item Every map with the right lifting property against $I$ is
      a trivial fibration: standard, levelwise.
  \end{enumerate}
  The two-out-of-three property for filtered quasi-isomorphisms
  follows from the long exact sequence of the filtration: a
  composite $V\to W\to U$ has filtered $E_1$-page agreeing with
  the composite of $E_1$-pages of $V\to W$ and $W\to U$ in each
  window, so two-out-of-three on each window propagates to the
  global filtered class. This proves the model structure in the
  chosen envelope only. No assertion is made for arbitrary unbounded
  Tate complexes outside it.
\end{proof}

\begin{rmk}[Boundedness restriction is harmless]
  Restriction to $\Cat{TateChain}_{\geq 0}^{\mathrm{adm}}$ avoids inverse-limit
  pathologies in unbounded filtrations. Every CE coalgebra
  $C_*^{\mathrm{CE}}(\mathfrak g)$ and every completed enveloping
  algebra $\widehat U(\mathfrak g)$ that appears here is
  non-negatively filtered (CE length and PBW length are
  $\geq 0$). A model structure on the raw unbounded category
  $\Cat{TateChain}$ would require the Spaltenstein semi-model
  modification (\cite{lurie-ha}*{\S 7.5}) and is not needed for
  the application here.
\end{rmk}

\subsubsection{Operad-algebra transfer}

Let $\mathrm{Lie}_{\Cat{TateVec}}$ be the Lie operad in
$\Cat{TateVec}$, regarded as a $\Sigma$-cofibrant $\Sigma$-module
(by the standard Koszul resolution of the $\Sigma$-cofibrant
replacement), and similarly $\mathrm{Ass}^+_{\Cat{TateVec}}$ for
augmented associative. A $\mathrm{Lie}_{\Cat{TateVec}}$-algebra is
a Tate Lie algebra in $\Cat{TateVec}$; an
$\mathrm{Ass}^+_{\Cat{TateVec}}$-algebra is a Tate complete
augmented associative algebra. Write
$$\mathrm{TateLie}=\mathrm{Lie}_{\Cat{TateVec}}\text{-alg}_{\geq 0},
  \qquad \mathrm{TateAssocAlg}^+=
  \mathrm{Ass}^+_{\Cat{TateVec}}\text{-alg}_{\geq 0}$$
in $\Cat{TateChain}_{\geq 0}^{\mathrm{adm}}$.

\begin{lem}[Transferred model structures on Tate operad-algebras]
  \label{lem:transferred-model}
  In the admissible envelope, the free-forget adjunction
  $\mathrm{Free}_{\mathcal P}\dashv\mathrm{Forget}$ for
  $\mathcal P\in\{\mathrm{Lie}_{\Cat{TateVec}},
  \mathrm{Ass}^+_{\Cat{TateVec}}\}$ transfers the cofibrantly
  generated model structure of \ref{thm:tate-model-structure} to
  $\mathcal P\text{-alg}_{\geq 0}$ in the sense of
  Hinich~\cite{hinich-htpy-alg}*{\S 2.2}: the weak equivalences and
  fibrations in $\mathcal P\text{-alg}_{\geq 0}$ are the maps
  whose underlying chain map in
  $\Cat{TateChain}_{\geq 0}^{\mathrm{adm}}$ is one,
  and the cofibrations are determined by lifting.
\end{lem}

\begin{proof}
  Hinich's transfer theorem requires:
  \begin{enumerate}[label=(\arabic*)]
    \item a path object on every algebra of the form
      $\mathrm{Free}_{\mathcal P}(F^w D^n)$;
    \item that pushouts along generating acyclic cofibrations be
      weak equivalences.
  \end{enumerate}
  Item~(1) follows because $\mathcal P=\mathrm{Lie}$ or
  $\mathrm{Ass}^+$ is a $\Sigma$-cofibrant operad in
  $\Cat{TateVec}$; the path object is constructed from the
  cylinder object on chain complexes via the Berger-Moerdijk
  path-functor lift~\cite{berger-moerdijk-axiomatic}, valid
  whenever the operad is admissible and the ambient category is
  cofibrantly generated and proper. Item~(2) follows because
  $\mathcal P$ acts levelwise on each Tate window, and the windowwise
  finite-dimensional pushout is a quasi-isomorphism by
  Schwede-Shipley \cite{schwede-shipley-monoid}*{Lemma 2.3}.
  Filtered-colimit assembly preserves quasi-isomorphism by
  Mittag-Leffler exactness (\alpharef\ item~1).
\end{proof}

The dual coalgebra side requires more care: Hinich's transfer
applies to algebra categories, not coalgebra categories, because
the cofree coalgebra functor is not a left adjoint in general.
The substitute is the Aubry-Lemaire-Smirnov \cite{loday-vallette}
\S 11.4 model structure on conilpotent dg coalgebras transferred
through the bar-cobar adjunction itself.

\begin{defn}[Conilpotent Tate coalgebras]\label{defn:tate-conilp-coalg}
  Let
  $$\begin{aligned}
  \mathrm{TateConilpCoAlg}
  =\bigl\{\,(C,\Delta_C,d_C)\,:\,&
  C\in\Cat{TateChain}_{\geq 0}^{\mathrm{adm}},\\
  &\Delta_C\colon C\to C\widehat\otimes C
    \text{ continuous coproduct},\\
  &\text{coaugmentation filtration exhaustive}\,\bigr\}.
  \end{aligned}$$
  with morphisms continuous coalgebra maps. Conilpotency means
  the coaugmentation filtration $F^{(n)}_{\mathrm{coaug}}C=
  \ker(\overline\Delta^{(n)})$ is exhaustive: $C=\bigcup_n
  F^{(n)}_{\mathrm{coaug}}C$.
\end{defn}

\begin{lem}[Transferred model structure on Tate conilpotent
  coalgebras]\label{lem:tate-conilp-model}
  $\mathrm{TateConilpCoAlg}$ admits, in the admissible envelope, a
  model structure with
  \begin{itemize}
    \item weak equivalences = maps $f\colon C\to C'$ such that
      $\Omega(f)\colon\Omega(C)\to\Omega(C')$ is a weak
      equivalence in $\mathrm{TateAssocAlg}^+$
      (\ref{lem:transferred-model});
    \item cofibrations = monomorphisms of underlying filtered
      chain complexes;
    \item fibrations = right lifting against acyclic cofibrations.
  \end{itemize}
  Equivalently, this is the Lefevre-Hasegawa
  \cite{lefevre-hasegawa}*{Th.~1.3.1.2} transferred model structure
  in the admissible Tate envelope.
\end{lem}

\begin{proof}
  Lefevre-Hasegawa, with proof sketched in Loday-Vallette
  \S 11.4.10, constructs the conilpotent-coalgebra transferred
  model structure in the symmetric monoidal category of $\Q$- or
  $\C$-vector spaces. The proof passes through the cobar-bar
  adjunction $\Omega\colon\mathrm{ConilpCoAlg}\rightleftarrows
  \mathrm{AugAssocAlg}\colon\mathrm{Bar}$ and uses that
  \begin{enumerate}[label=(\roman*)]
    \item $\Omega$ takes coalgebra cofibrations (= mono with
      filtered cokernel) to algebra cofibrations;
    \item $\Omega$ inverts coalgebra weak equivalences;
    \item the unit $\eta_C\colon C\to\mathrm{Bar}\,\Omega(C)$ is a
      weak equivalence for every conilpotent dg coalgebra.
  \end{enumerate}
  Items~(i)-(iii) are checked windowwise (each Tate window is
  finite-dimensional, so the standard Lefevre-Hasegawa proof
  applies) and assemble by Mittag-Leffler exactness. The unit
  weak-equivalence~(iii) is precisely \alpharef\ item~3
  (admissible filtered qiso). Therefore the construction lifts
  verbatim to the admissible Mittag-Leffler Tate envelope.
\end{proof}

\subsubsection{Bar-cobar Quillen equivalence in the admissible envelope}

\begin{thm}[Tate bar-cobar Quillen equivalence in the admissible envelope]
  \label{thm:tate-bar-cobar-quillen}
  The bar-cobar adjunction
  $$\Omega\colon\mathrm{TateConilpCoAlg}\rightleftarrows
    \mathrm{TateAssocAlg}^+\colon\mathrm{Bar}$$
  is a Quillen equivalence with respect to the admissible-envelope
  model structures
  of \ref{lem:transferred-model} and \ref{lem:tate-conilp-model}.
  In particular:
  \begin{enumerate}
    \item $\Omega$ is left Quillen: it preserves cofibrations and
      acyclic cofibrations.
    \item $\mathrm{Bar}$ is right Quillen: it preserves
      fibrations and acyclic fibrations.
    \item For every cofibrant conilpotent coalgebra $C$ and every
      fibrant augmented associative algebra $A$, the unit
      $\eta_C\colon C\to\mathrm{Bar}\,\Omega(C)$ and the counit
      $\varepsilon_A\colon\Omega\,\mathrm{Bar}(A)\to A$ are weak
      equivalences.
    \item The induced derived adjunction
      $L\Omega\colon\mathrm{Ho}(\mathrm{TateConilpCoAlg})
      \rightleftarrows\mathrm{Ho}(\mathrm{TateAssocAlg}^+)\colon
      R\,\mathrm{Bar}$ is an equivalence of homotopy categories.
  \end{enumerate}
\end{thm}

\begin{proof}
  Items~(1) and~(2) are immediate from
  \ref{lem:tate-conilp-model}: weak equivalences in
  $\mathrm{TateConilpCoAlg}$ are defined as maps whose $\Omega$
  is a weak equivalence, and fibrations on the coalgebra side are
  defined by lifting.

  For item~(3), the unit $\eta_C$: filter both sides by symmetric
  / tensor length on the underlying $\Cat{TateVec}$ object, plus
  the Tate window filtration $F^w$. The associated graded gives,
  windowwise, a finite-dimensional dg coalgebra and the standard
  bar-cobar resolution
  $C^{\mathrm{gr}}\to\mathrm{Bar}\,\Omega(C^{\mathrm{gr}})$. By
  Loday-Vallette \cite{loday-vallette}*{\S 2.2.5, \S 11.1.5}, this
  is a quasi-isomorphism on each window. Conilpotency
  (\ref{defn:tate-conilp-coalg}) ensures the coaugmentation
  filtration on $C$ is exhaustive and Hausdorff, so the abutment
  of the spectral sequence of the filtration computes the global
  qiso. Mittag-Leffler exactness in the Tate filtration
  (\alpharef\ item~1) lifts the windowwise qiso to a global
  filtered qiso.

  For the counit $\varepsilon_A$: filter $\widehat U(\mathfrak g)$
  by PBW. \alpharef\ item~4 gives
  $\mathrm{gr}_F^\bullet\widehat U(\mathfrak g)
   \cong\widehat S_{\mathrm{Tate}}(\mathfrak g)$
  (PBW completed-symmetric associated graded), and item~3 gives
  the windowwise bar-cobar acyclicity for the Lie operad. The
  symmetric Koszul resolution of Loday-Vallette
  \cite{loday-vallette}*{\S 3.4.13} is qiso on associated graded.
  The same Mittag-Leffler argument propagates to the abutment.

  Item~(4) is formal from~(1)-(3) by Quillen's adjoint-functor
  theorem (\cite{hovey}*{Th.~1.3.10}).
\end{proof}

\subsubsection[Promotion to the associated infinity-categories]{Promotion to the associated $\infty$-categories}

\begin{cor}[$\infty$-categorical Tate bar-cobar equivalence]
  \label{cor:tate-bar-cobar-infinity}
  Apply the Dwyer-Kan / Hammock-localization functor to the
  Quillen equivalence of \ref{thm:tate-bar-cobar-quillen}:
  $$N_\Delta\bigl(\mathrm{TateConilpCoAlg}^{\mathrm{cf}},W\bigr)
    \simeq
    N_\Delta\bigl(\mathrm{TateAssocAlg}^{+,\mathrm{cf}},W\bigr)$$
  as the $\infty$-categories associated to the admissible model
  envelope in the sense of Lurie's
  \emph{Higher Algebra} \S A.3, where $\mathrm{cf}$ denotes the
  full subcategory of cofibrant-fibrant objects and $W$ is the
  class of weak equivalences. The equivalence is induced by
  $L\Omega$ and $R\,\mathrm{Bar}$ at the cofibrant-fibrant level.
  Presentability is an attribute of the chosen locally presentable
  envelope, not of the raw Tate category.
\end{cor}

\begin{proof}
  Lurie \cite{lurie-ha}*{A.3.7.6, A.3.7.10}: a Quillen equivalence
  of model categories induces an equivalence of associated
  $\infty$-categories. The presentability of both sides, inside the
  envelope, is
  inherited from the cofibrant generation of the model structures
  (\ref{thm:tate-model-structure}, \ref{lem:transferred-model},
  \ref{lem:tate-conilp-model}) by Lurie A.2.6.8.
\end{proof}

\subsubsection{Identification with the Lie operadic Koszul duality}

The Quillen equivalence of \ref{thm:tate-bar-cobar-quillen} is
\emph{associative} bar-cobar. Specialization to the Lie operad
gives the cochain-level equivalence used in the CE/PV theorem lane.

\begin{cor}[Tate Lie bar-cobar / CE Koszul duality]
  \label{cor:tate-lie-bar-cobar}
  For $\mathfrak g$ a conilpotent Tate graded Lie algebra in the
  sense of \alpharef, the bar-cobar Quillen equivalence applied
  to $A=\widehat U(\mathfrak g)\in\mathrm{TateAssocAlg}^+$ yields
  the equivalence in the homotopy category
  $$L\Omega\bigl(C_*^{\mathrm{CE}}(\mathfrak g)\bigr)
    \xrightarrow{\sim}\widehat U(\mathfrak g),\qquad
    R\,\mathrm{Bar}\bigl(\widehat U(\mathfrak g)\bigr)
    \xrightarrow{\sim}C_*^{\mathrm{CE}}(\mathfrak g),$$
  and hence a Lie-Koszul-dual equivalence of cochain complexes
  $$C^\bullet_{\mathrm{CE},\mathrm{cont}}(\mathfrak g)
    \simeq\bigl(\widehat U(\mathfrak g)\bigr)^!$$
  in $\mathrm{Ho}(\mathrm{TateConilpCoAlg})$.

  Inside the admissible envelope, by the symmetric-monoidal-functoriality of the bar-cobar
  Quillen equivalence \cite{loday-vallette}*{\S 13.2} and Lurie
  \cite{lurie-ha}*{\S 5.5}, this lifts to an equivalence of
  symmetric monoidal $\infty$-categories,
  $$N_\Delta\bigl(\mathrm{TateConilpCoAlg}_{\mathrm{Lie}}^{\mathrm{cf}},
    W\bigr)\simeq
    N_\Delta\bigl(\mathrm{TateAssocAlg}_{\mathrm{Lie}}^{+,\mathrm{cf}},
    W\bigr)$$
  in the framework of Lurie \emph{Higher Algebra} \S 5.5
  (locally constant factorization algebras over $\R$, applied
  pointwise on the line; or directly the formal-disk version).
\end{cor}

\begin{proof}
  Apply \ref{thm:tate-bar-cobar-quillen} to the
  Lie-operadic Koszul-dual pair $(\mathrm{Lie},\mathrm{CoComm})$
  in $\Cat{TateVec}$. Loday-Vallette \cite{loday-vallette}*{Th.~13.4.6}
  identifies the operadic bar-cobar with the Lie-bar-cobar of
  CE chains and completed enveloping algebra. The
  symmetric-monoidal extension is direct because the operad
  $\mathrm{Lie}$ is a symmetric-monoidal Hopf operad in
  $\Cat{Vec}$, and base change to $\Cat{TateVec}$ via the strong
  symmetric-monoidal functor $\Cat{Vec}\to\Cat{TateVec}$
  preserves operadic Koszul duality
  (\cite{loday-vallette}*{\S 13.5}).
\end{proof}

\subsubsection{Application to the local theorem lanes}

\begin{thm}[Admissible-envelope promotion of the CE/PV cochain identity]
\label{thm:open-closed-derived-center-derived}
  Let $\mathfrak h_{\mathrm{adm}}$ be a pro-finite-dimensional Lie algebra in
  $\Cat{TateVec}$ whose shifted-cotangent extension
  $\mathfrak g_{\mathrm{adm}}
  =\mathfrak h_{\mathrm{adm}}\ltimes
  \mathfrak h_{\mathrm{adm},\mathrm{cont}}^\vee[1]$
  satisfies the Tate-conilpotent admissible-envelope hypotheses of
  \ref{thm:tate-bar-cobar-quillen}.
  The strict cochain map
  $$\Phi\colon C^\bullet_{\mathrm{CE},\mathrm{cont}}(\mathfrak g_{\mathrm{adm}})
    \longrightarrow \mathrm{PV}\bigl(
    \widehat S_{\mathrm{Tate}}(\mathfrak h_{\mathrm{adm}})\bigr)$$
  of \ref{thm:open-closed-derived-center} promotes to an
  equivalence, inside that admissible Tate model envelope, of module
  $\infty$-categories over the two identified completed dg
  \(P_0\)-algebras
  $$\widetilde\Phi\colon
    N_\Delta\bigl(C^\bullet_{\mathrm{CE},\mathrm{cont}}(\mathfrak g_{\mathrm{adm}})
      \text{-mod},W\bigr)
    \xrightarrow{\sim}
    N_\Delta\bigl(\mathrm{PV}\bigl(\widehat S_{\mathrm{Tate}}(\mathfrak h_{\mathrm{adm}})
      \bigr)\text{-mod},W\bigr)$$
  in the framework of Lurie's \emph{Higher Algebra} \S 5.5.  This is
  not a construction of a completed dg \(P_0\) model structure or a
  factorization-centre theorem for the full untruncated formal Hamiltonian
  algebra. The equivalence itself is the admissible-envelope algebraic
  one.
\end{thm}

\begin{proof}
  By \ref{cor:tate-lie-bar-cobar}, the cochain-level identity
  $$\Phi\colon C^\bullet_{\mathrm{CE},\mathrm{cont}}(\mathfrak g_{\mathrm{adm}})
    \xrightarrow{\sim}
    \widehat S_{\mathrm{Tate}}(\mathfrak h_{\mathrm{adm}})\widehat\otimes
    \widehat\Lambda(\theta^I)$$
  of \ref{thm:open-closed-derived-center} (which strictly
  intertwines $d_{\mathrm{CE}}$ and $d_\pi$ by
  \ref{thm:open-closed-derived-center}) is the underlying $0$-cell
  of an equivalence in
  $N_\Delta(\mathrm{TateConilpCoAlg}_{\mathrm{Lie}}^{\mathrm{cf}})$.
  The right-hand side, $\mathrm{PV}\bigl(\widehat S_{\mathrm{Tate}}
  (\mathfrak h_{\mathrm{adm}})\bigr)$, has the standard Schouten-Lichnerowicz
  $P_0$-algebra structure (\ref{lem:linear-poisson-schouten}),
  and the closed-side cochain complex inherits a $P_0$-algebra
  structure by transport via $\Phi$.

  The $\infty$-category structure on completed dg
  $P_0$-algebras follows from
  \ref{thm:tate-bar-cobar-quillen} applied to the operad
  $P_0=\mathrm{Comm}\circ\Lambda^\bullet[\,\cdot\,,\,\cdot\,]_S$,
  which is Koszul self-dual up to suspension by Cattaneo-Felder
  \cite{cattaneo-felder-bv} (cf. Costello-Gwilliam Vol.~II
  \S 1.4 for the Tate version).

  The equivalence on the module category follows from
  \cite{lurie-ha}*{Th.~7.1.4.6}: a Quillen equivalence of
  $E_n$-monoidal model categories induces an equivalence of
  module $\infty$-categories. Applied with $n=1$ (associative
  monoidal in the bar-cobar; $P_0$ structure superimposed by
  the Lichnerowicz-Schouten differential) gives
  $\widetilde\Phi$.

  Functoriality in the Lie algebra gives interval-wise compatibility
  for admissible interval objects
  $\mathfrak h_{\mathrm{adm},I}=
  \Omega^\bullet_c(I)\otimes\mathfrak h_{\mathrm{adm}}$ when these
  interval objects remain in the same admissible envelope.  This is
  interval-wise categorical compatibility, not an independent Weiss
  descent theorem and not an unreduced open-BV centrality theorem.
\end{proof}

\begin{rmk}[Relation to Lurie's higher-algebraic Koszul
  duality]
  \label{rmk:lurie-koszul}
  Lurie \cite{lurie-ha}*{\S 5.5.5} constructs a higher-algebraic
  Koszul duality between augmented $\infty$-operad-algebras and
  their Koszul-dual coalgebras. For the operad
  $\mathrm{Lie}_{\Cat{TateVec}}$, Koszul self-duality (with
  shift) gives the equivalence of \ref{cor:tate-lie-bar-cobar}
  in the Tate-conilpotent setting. The advantage of the model-
  categorical proof is concrete: weak equivalences are spelled
  out in terms of admissible filtered qiso of \alpharef, not
  abstractly via the Hammock localization. The two constructions
  agree by Lurie A.3.7 plus the Hinich transfer of
  \ref{lem:transferred-model}.
\end{rmk}

\subsubsection{Compatibility with the descendant residual}

The Tate Quillen equivalence of \ref{thm:tate-bar-cobar-quillen}
operates on the bar-cobar duality between $\widehat U(\mathfrak g)$
and $C_*^{\mathrm{CE}}(\mathfrak g)$. It does not directly use
the Casimir element
$C_{\mathfrak h}=\sum_d e_{d,i}\otimes e_d^{i,\vee}$
of \alpharef\ or its convergence in the Tate completion.

\begin{prop}[Quillen equivalence vs Casimir convergence]
  \label{prop:quillen-vs-casimir}
  The Quillen equivalence of \ref{thm:tate-bar-cobar-quillen}
  and the resulting $\infty$-categorical equivalence of
  \ref{thm:open-closed-derived-center-derived}, when their
  admissible-envelope hypotheses hold, are independent
  of the convergence of the Tate Casimir element. The cochain
  complexes $C^\bullet_{\mathrm{CE},\mathrm{cont}}(\mathfrak g)$
  and $\mathrm{PV}\bigl(\widehat S_{\mathrm{Tate}}(\mathfrak h)
  \bigr)$ are equivalent in the model category of
  \ref{lem:transferred-model} only after the same conilpotent
  admissibility hypotheses are imposed. Casimir convergence is not one
  of those hypotheses; conilpotence is.
\end{prop}

\begin{proof}
  The bar-cobar comparison uses
  $\widehat S^c(\mathfrak g[1])$ and $T(\mathfrak g)$ as
  underlying graded objects. Both are well-defined and conilpotent
  in the admissible envelope under \alpharef\ items~1, 2 (continuous
  dual exact, conilpotent CE coalgebra). The Casimir
  $\sum_d e_{d,i}\otimes e_d^{i,\vee}$ is an element of
  $\mathfrak h\widehat\otimes\mathfrak h^\vee_{\mathrm{cont}}$; it
  is required for the Costello \cite{costello-renormalization}
  heat-kernel propagator construction at the analytic /
  factorization-algebra-with-non-strict-bracket level. It is
  \emph{not} required for the algebraic bar-cobar comparison.

  Therefore the cochain-level $\Phi$ promotes, in the admissible
  model envelope and only there, to an $\infty$-categorical equivalence. What
  remains outside this algebraic promotion is
  the Weiss-cosheaf upgrade and the heat-kernel quantum lift,
  flagged in \ref{prop:weiss-cosheaf-residual} below.
\end{proof}

\begin{prop}[Residual Weiss-cosheaf condition]
  \label{prop:weiss-cosheaf-residual}
  The $\infty$-categorical equivalence of
  \ref{thm:open-closed-derived-center-derived} on stalks and on
  the locally-constant interval-factorization $\mathcal F_{\mathrm{cl}}$
  of \ref{thm:open-closed-derived-center-factorization} does not,
  by itself, give the Weiss-cosheaf condition for the
  \emph{unreduced} factorization algebra
  $\mathrm{Obs}^{\mathrm{cl}}_{\mathrm{closed},H}$ on $\R^2\times\C^2$.
  The Weiss / descent condition for an unreduced factorization
  algebra on a higher-dimensional manifold requires checking
  \v{C}ech descent for Weiss covers, which is a separate
  cosheaf-theoretic statement.  Costello-Gwilliam Vol.~I
  \S\S 6.1, 6.5, 7.2 provide the published factorization,
  cosheaf, and basis-extension vocabulary; no theorem from those
  sections is used here to close the unreduced descent obligation.  T3
  closes the cochain-level $\infty$-categorical promotion; the Weiss
  upgrade is recorded as the standing residual on the
  unreduced-factorization side of the local theorem.
\end{prop}

\begin{proof}
  Lurie \cite{lurie-ha}*{\S 5.5.4}, with the same terminology as
  Costello-Gwilliam \cite{costello-gwilliam}*{Vol.~I \S 6.4}: a locally constant factorization
  algebra on $\R$ determines an $E_1$-algebra up to the standard
  contractible choice; the
  Quillen equivalence of \ref{thm:tate-bar-cobar-quillen}
  realizes this. On a higher-dimensional manifold, the Weiss
  cover $\{U_\alpha\}$ has structure-map descent independent of
  the algebraic bar-cobar; in particular the heat-kernel
  propagator on $\R^2\times\C^2$ enters the analytic check.
  Without that propagator the bar-cobar identity is consistent
  but does not trivialize Weiss descent.
\end{proof}

\subsubsection{Summary diagram}

\begin{equation*}
  \begin{array}{ccc}
    \alpharef\text{ (filtered qiso of cochain complexes)} &
    \xrightarrow{\text{T3 promotion}} &
    \text{Quillen equivalence of model cats} \\
    \downarrow & & \downarrow \\
    \text{cochain identity }\Phi\,d_{\mathrm{CE}}=d_\pi\,\Phi &
    \xrightarrow{\text{Lurie A.3.7}} &
    \text{equivalence of }\infty\text{-cats}\\
    & & \downarrow\text{Lurie 5.5} \\
    & &
	    \text{equivalence in admissible }\infty\text{-cats of}\\
    & & \text{completed dg }P_0\text{-algebras}
  \end{array}
\end{equation*}

The chain on the left is $\Phi$ as a strict cochain map.
The chain on the right is T3's promotion: model-categorical,
then $\infty$-categorical, then $P_0$-algebra-categorical.
The promotion does not depend on the Tate Casimir convergence
(\ref{prop:quillen-vs-casimir}); the Weiss-cosheaf descent on the
unreduced factorization algebra is a separate residual
(\ref{prop:weiss-cosheaf-residual}).
